% Volume V: Dynamics, Signals, Graphs, and Control
% Continuity note for Chapter 26.
% Previous dependencies: Chapter 8 introduced derivatives, integrals, and ODEs; Chapter 9 introduced partial derivatives, gradients, Jacobians, and Hessians; Chapter 12 introduced function spaces, operators, Fourier modes, sampling, and filters; Chapters 24-25 introduced discretization, numerical stability, ODE flows, linearization, and gradient dynamics.
% New durable concepts: fields over space and time; partial differential equations; initial and boundary conditions; gradient, divergence, and Laplacian operators; conservation laws; advection, diffusion, reaction, heat, and wave equations; Fourier-mode solutions; finite-difference discretization; reaction-diffusion pattern intuition; neural-field equations; continuous-depth models, neural ODEs, adjoint sensitivities, solver tolerances, stiffness, and scientific machine learning.
% Forward hooks: Chapter 27 adds stochastic processes and SDEs; Chapter 28 studies spectral analysis of signals; Chapter 30 uses dynamics in control; Chapters 43 and 49 use cable equations and neural fields; Chapters 40-42 use continuous-depth and generative dynamics in deep learning.
% Notation commitments: Spatial variables are x in \Omega subset R^d; time is t; fields are u(x,t) or vector fields u:\Omega x [0,T] -> R^m; partial derivatives are u_t=\partial u/\partial t and u_{xx}=\partial^2u/\partial x^2; gradient, divergence, and Laplacian are \nabla u, \nabla\cdot F, and \Delta u; boundary conditions are Dirichlet, Neumann, or periodic.
\section{Chapter 26: Partial Differential Equations and Continuum Models}
\label{ch:partial-differential-equations-continuum-models}

Chapter 25 studied states \(x(t)\in\mathbb{R}^n\) that evolve over time. Many scientific systems need one more independent variable: space. A membrane voltage can vary along a dendrite, an activity profile can vary across cortex, a concentration can vary across tissue, and an image can be treated as a function over a two-dimensional domain. Partial differential equations are the language for fields that vary over space and time.

The new issue is that a field has infinitely many degrees of freedom before discretization. To determine a solution, we need not only a local differential law but also initial conditions, boundary conditions, and a clear spatial domain. Numerical computation then turns the continuum model into a finite-dimensional approximation.

\paragraph{Chapter problem.}
How do local rules involving time and spatial derivatives determine the evolution of fields, and how do diffusion, waves, reactions, neural fields, and continuous-depth models use this continuum language?

\paragraph{Section map.}
Section 26.1 defines PDEs, fields, spatial operators, and boundary conditions. Section 26.2 develops diffusion, heat, advection, and wave equations. Section 26.3 studies reaction-diffusion systems and neural fields. Section 26.4 connects continuous-depth models, neural ODEs, adjoint gradients, and scientific machine learning.

\subsection{26.1 PDEs and fields}

\paragraph{Guiding question.}
What changes mathematically when the unknown is a field \(u(x,t)\) rather than a finite-dimensional state \(x(t)\)?

\paragraph{Beginner intuition.}
An ODE tracks a state over time. A PDE tracks a field over space and time. For an ODE, the unknown might be a voltage \(V(t)\). For a PDE, the unknown might be a voltage profile \(V(x,t)\), where \(x\) is position along a cable. The derivative \(dV/dt\) is replaced by partial derivatives such as
\[
\frac{\partial V}{\partial t},
\qquad
\frac{\partial V}{\partial x},
\qquad
\frac{\partial^2 V}{\partial x^2}.
\]
A partial derivative changes one independent variable while holding the others fixed.

\paragraph{Formal setup.}
Let \(\Omega\subseteq\mathbb{R}^d\) be a spatial domain and let \(t\in[0,T]\). A scalar field is a function
\[
u:\Omega\times[0,T]\to\mathbb{R},
\qquad
(x,t)\mapsto u(x,t).
\]
A vector-valued field has codomain \(\mathbb{R}^m\). We write
\[
u_t=\frac{\partial u}{\partial t},
\qquad
u_{x_i}=\frac{\partial u}{\partial x_i}.
\]
The spatial gradient is
\[
\nabla u=\begin{bmatrix}u_{x_1}\\ \vdots\\ u_{x_d}\end{bmatrix},
\]
the divergence of a vector field \(F:\Omega\to\mathbb{R}^d\) is
\[
\nabla\cdot F=\sum_{i=1}^d\frac{\partial F_i}{\partial x_i},
\]
and the Laplacian is
\[
\Delta u=\nabla\cdot\nabla u=\sum_{i=1}^d\frac{\partial^2u}{\partial x_i^2}.
\]

A PDE is an equation involving \(u\) and some of its partial derivatives, such as
\[
u_t=D\Delta u.
\]
To determine a solution, one usually specifies an initial condition
\[
u(x,0)=u_0(x)
\]
and boundary conditions on \(\partial\Omega\). Common boundary conditions include Dirichlet,
\[
u(x,t)=g(x,t)\quad\text{on }\partial\Omega,
\]
Neumann,
\[
\nabla u(x,t)\cdot n(x)=h(x,t)\quad\text{on }\partial\Omega,
\]
and periodic boundary conditions, where opposite sides of the domain are identified.

\paragraph{Worked mechanism: advection as transported shape.}
The one-dimensional advection equation is
\[
u_t+c u_x=0,
\]
where \(c\) is a constant velocity. If the initial field is \(u(x,0)=u_0(x)\), then
\[
\boxed{u(x,t)=u_0(x-ct)}
\]
solves the PDE. To check this, let \(s=x-ct\). Then
\[
u_t=-c u_0'(s),\qquad u_x=u_0'(s),
\]
so
\[
u_t+c u_x=-c u_0'(s)+c u_0'(s)=0.
\]
The field profile is transported to the right when \(c>0\) without changing shape.

\paragraph{Examples and connections.}
A numerical example is \(u_0(x)=\sin x\), \(c=2\). Then
\[
u(x,t)=\sin(x-2t).
\]
At time \(t=1\), the whole sine wave has shifted by \(2\) spatial units.

In computational neuroscience, a dendritic cable voltage \(V(x,t)\) is a field over position \(x\) along the dendrite and time \(t\). Boundary conditions encode what happens at the soma, sealed end, or coupled branch point. Treating all dendritic locations as one compartment would turn the model into an ODE, but the PDE keeps spatial voltage gradients.

In AI and scientific machine learning, images, physical states, and spatial latent fields are often represented as functions on a domain. Neural operators and physics-informed neural networks learn maps between functions or solve PDE-constrained problems; their inputs and outputs are fields, not just vectors.

\paragraph{Common misconceptions and failure modes.}
A PDE is not just a harder ODE. The solution depends on spatial domain and boundary conditions. A partial derivative is not a total derivative; \(u_t\) holds space fixed while time changes. Another failure is to discretize space without checking whether the boundary condition was discretized consistently. Finally, a field may be vector-valued, so one PDE system can contain several coupled scalar equations.

\paragraph{Exercises.}
\begin{enumerate}[label=\textbf{26.1.\arabic*.}, leftmargin=*]
\item \textbf{Mechanical.} For \(u(x,t)=x^2t+\sin t\), compute \(u_t\), \(u_x\), and \(u_{xx}\).
\item \textbf{Proof or derivation.} Verify that \(u(x,t)=u_0(x-ct)\) solves \(u_t+c u_x=0\) when \(u_0\) is differentiable.
\item \textbf{Numerical/computational.} On grid points \(x_i=ih\), write the centered finite-difference approximation for \(u_x(x_i,t)\) using \(u_{i+1}\) and \(u_{i-1}\).
\item \textbf{Interpretation.} What does a Neumann boundary condition with \(\nabla u\cdot n=0\) mean for flux through a boundary?
\item \textbf{Misconception/debugging.} A student gives a heat equation and initial condition but no boundary condition on a finite interval. Explain why the model is incomplete.
\end{enumerate}

\subsection{26.2 Diffusion, heat, and wave equations}

\paragraph{Guiding question.}
How do canonical PDEs represent smoothing, transport, and propagation?

\paragraph{Beginner intuition.}
Different spatial derivatives encode different mechanisms. A first derivative often appears in transport or advection: a profile moves. A Laplacian often appears in diffusion: high regions spread into low regions and the field smooths out. A second time derivative appears in waves: disturbances propagate and oscillate rather than simply smooth away.

\paragraph{Formal setup.}
The diffusion or heat equation is
\[
u_t=D\Delta u,\qquad D>0.
\]
In one dimension,
\[
u_t=D u_{xx}.
\]
The advection-diffusion equation combines transport and smoothing:
\[
u_t+c u_x=D u_{xx}.
\]
The wave equation is
\[
u_{tt}=c^2\Delta u,
\]
where \(c>0\) is wave speed.

Diffusion can be derived from conservation. Let \(u(x,t)\) be a density and \(J(x,t)\) a flux. Conservation in one dimension says
\[
u_t+J_x=0.
\]
Fick's law says flux moves down gradients:
\[
J=-D u_x.
\]
Substituting gives
\[
u_t-Du_{xx}=0,
\]
or \(u_t=D u_{xx}\).

\paragraph{Worked mechanism: Fourier modes for heat and waves.}
On a periodic one-dimensional domain, consider a Fourier mode
\[
u(x,t)=a(t)\sin(kx).
\]
For the heat equation \(u_t=D u_{xx}\),
\[
u_t=a'(t)\sin(kx),
\qquad
u_{xx}=-k^2a(t)\sin(kx).
\]
Thus
\[
a'(t)=-Dk^2a(t),
\]
so
\[
a(t)=a(0)e^{-Dk^2t}.
\]
High spatial frequencies, with large \(k\), decay faster. This is the mathematical reason diffusion smooths.

For the wave equation, the same spatial mode gives
\[
a''(t)=-c^2k^2a(t),
\]
so
\[
a(t)=A\cos(ckt)+B\sin(ckt).
\]
The mode oscillates instead of decaying.

\paragraph{Examples and connections.}
A concrete heat example is
\[
u(x,0)=\sin(3x)
\]
on a periodic domain. The solution is
\[
u(x,t)=e^{-9Dt}\sin(3x).
\]
If \(D=0.1\), the amplitude at \(t=1\) is \(e^{-0.9}\) times the initial amplitude.

In computational neuroscience, diffusion-like coupling appears in dendritic cable equations and spatially extended activity models. Voltage differences along a cable drive axial current, producing a Laplacian-like spatial smoothing term. Waves appear in models of propagating activity, but a wave model is an idealization whose validity depends on the tissue, variables, and scale.

In AI, diffusion-like PDEs appear in image smoothing, score-based generative modeling analogies, and scientific ML surrogates. The PDE heat equation is deterministic smoothing; stochastic diffusion models add probabilistic noise and reverse-time learning, which Chapter 27 and Chapter 41 will distinguish carefully.

\paragraph{Common misconceptions and failure modes.}
Diffusion is not the same as advection. Diffusion smooths and spreads; advection transports. The heat equation and wave equation may both use spatial derivatives, but their time behavior is qualitatively different. Another failure is to forget boundary conditions: heat on a finite interval with fixed endpoints behaves differently from heat on a periodic circle. Finally, a finite-difference discretization of diffusion has a time-step stability constraint for explicit methods.

\paragraph{Exercises.}
\begin{enumerate}[label=\textbf{26.2.\arabic*.}, leftmargin=*]
\item \textbf{Mechanical.} For \(u(x,t)=e^{-4Dt}\sin(2x)\), verify that \(u_t=D u_{xx}\).
\item \textbf{Proof or derivation.} Derive the heat equation from conservation \(u_t+J_x=0\) and Fick's law \(J=-D u_x\).
\item \textbf{Numerical/computational.} Approximate \(u_{xx}(x_i,t)\) by \((u_{i+1}-2u_i+u_{i-1})/h^2\). If \(u_{i-1}=1\), \(u_i=3\), and \(u_{i+1}=2\), compute the approximation for \(h=0.5\).
\item \textbf{Interpretation.} Why do high-frequency Fourier modes decay faster under the heat equation?
\item \textbf{Misconception/debugging.} A student says, "The wave equation smooths out bumps like the heat equation." Use the Fourier-mode calculation to correct the statement.
\end{enumerate}

\subsection{26.3 Reaction-diffusion and neural fields}

\paragraph{Guiding question.}
How can local nonlinear dynamics and spatial coupling create patterns, bumps, and waves?

\paragraph{Beginner intuition.}
Diffusion alone smooths. Reaction alone changes each point independently according to local ODE dynamics. Reaction-diffusion systems combine the two:
\[
u_t=D\Delta u+R(u).
\]
The term \(D\Delta u\) couples neighboring spatial locations. The reaction term \(R(u)\) supplies local growth, saturation, inhibition, excitation, or chemical kinetics. Together they can create spatial patterns that neither part would create alone.

\paragraph{Formal setup.}
A scalar reaction-diffusion equation is
\[
u_t=D\Delta u+R(u),
\]
where \(D>0\) and \(R:\mathbb{R}\to\mathbb{R}\). A vector reaction-diffusion system has
\[
u_t=D\Delta u+R(u),
\qquad
u(x,t)\in\mathbb{R}^m,
\]
where \(D\) may be a diagonal matrix of diffusion constants and \(R:\mathbb{R}^m\to\mathbb{R}^m\).

A neural-field model often uses an integral spatial coupling instead of a Laplacian:
\[
\tau u_t(x,t)
=-u(x,t)+\int_\Omega w(x-y)S(u(y,t))\,dy+I(x,t).
\]
Here \(u(x,t)\) is activity at position \(x\), \(w(x-y)\) is a spatial connectivity kernel, \(S\) is a static nonlinearity, and \(I\) is input. The integral sums activity from other locations weighted by distance or feature similarity.

\paragraph{Worked mechanism: linear stability of a homogeneous state.}
Suppose \(u^*\) is a spatially homogeneous equilibrium of
\[
u_t=D\Delta u+R(u),
\]
meaning \(R(u^*)=0\). Let \(u(x,t)=u^*+\eta(x,t)\) with small perturbation \(\eta\). Linearizing gives
\[
\eta_t=D\Delta \eta+R'(u^*)\eta.
\]
For a Fourier mode \(\eta(x,t)=a(t)\sin(kx)\),
\[
a'(t)=\left(R'(u^*)-Dk^2\right)a(t).
\]
The mode grows if
\[
R'(u^*)-Dk^2>0
\]
and decays otherwise. In scalar systems diffusion suppresses high-frequency growth. In multi-component systems with different diffusion constants and cross-coupled reactions, diffusion can destabilize some spatial modes, giving pattern formation. That is the beginning of Turing instability intuition.

\paragraph{Examples and connections.}
A scalar example is the Fisher-KPP equation
\[
u_t=D u_{xx}+r u(1-u),
\]
where \(0\leq u\leq 1\). The reaction term grows small positive \(u\) and saturates near \(1\), while diffusion spreads the activity. This equation supports traveling fronts.

In computational neuroscience, a bump attractor can be modeled by a neural field in which nearby positions excite and distant positions inhibit. A localized input can create a stable activity bump representing a remembered spatial location, head direction, or continuous stimulus feature. The model is a continuum approximation to many neurons arranged along a feature coordinate.

In AI and scientific ML, reaction-diffusion equations are used as benchmark systems for learned simulators, neural operators, and PDE-constrained learning. Convolutional networks can be interpreted loosely as local spatial coupling plus pointwise nonlinear reaction, but the analogy must be used carefully because standard CNN layers are discrete learned maps, not necessarily time-discretized PDEs.

\paragraph{Common misconceptions and failure modes.}
Reaction-diffusion does not mean every pattern in biology is explained by a Turing mechanism. The model assumptions and parameters matter. A neural field is also not a literal map of every neuron; it is a coarse-grained activity model over a spatial or feature coordinate. Another failure is to ignore the kernel \(w\): local excitation, global inhibition, and oscillatory kernels produce different dynamics. Finally, linear stability predicts small-perturbation behavior, not the final nonlinear pattern by itself.

\paragraph{Exercises.}
\begin{enumerate}[label=\textbf{26.3.\arabic*.}, leftmargin=*]
\item \textbf{Mechanical.} For \(R(u)=u(1-u)\), find the homogeneous equilibria of \(u_t=D\Delta u+R(u)\).
\item \textbf{Proof or derivation.} Linearize \(u_t=D\Delta u+R(u)\) around a homogeneous equilibrium \(u^*\) and derive \(a'(t)=(R'(u^*)-Dk^2)a(t)\) for a sine Fourier mode.
\item \textbf{Numerical/computational.} In the neural-field equation, suppose \(w(x-y)\) is nonzero only for nearby positions. What sparsity or locality would this create after discretizing space?
\item \textbf{Interpretation.} In a bump attractor model, what does the location of the activity bump represent?
\item \textbf{Misconception/debugging.} A student says, "Diffusion always destroys patterns, so reaction-diffusion cannot create them." Explain why multi-component reaction-diffusion systems can destabilize some modes.
\end{enumerate}

\subsection{26.4 Continuous-depth models}

\paragraph{Guiding question.}
How do deep learning models use continuous-time dynamics, and what numerical issues does that introduce?

\paragraph{Beginner intuition.}
A residual network updates a hidden state by adding a learned increment:
\[
h_{k+1}=h_k+\Delta t\, f_\theta(h_k,t_k).
\]
This looks like an Euler step for the ODE
\[
\dot h(t)=f_\theta(h(t),t).
\]
A continuous-depth model treats depth as time and computes an output by solving an ODE from an initial hidden state to a final hidden state. The model is not defined only by \(f_\theta\); it is also defined by the numerical solver, tolerances, and gradient method used during training.

\paragraph{Formal setup.}
A neural ODE has the form
\[
\frac{dh}{dt}=f_\theta(h(t),t),
\qquad h(t_0)=h_0,
\]
and output
\[
h(t_1)=\Phi_{t_0,t_1}^{\theta}(h_0),
\]
where \(\Phi\) is the flow map determined by the learned vector field and the solver. If a loss depends on the final state,
\[
L=\ell(h(t_1)),
\]
then gradients can be computed by differentiating through the solver or by an adjoint equation. In a continuous adjoint derivation, define
\[
a(t)=\frac{\partial L}{\partial h(t)}.
\]
Under smoothness assumptions,
\[
\frac{da}{dt}=-\left(\frac{\partial f_\theta}{\partial h}(h(t),t)\right)^\top a(t),
\]
integrated backward from
\[
a(t_1)=\nabla_h \ell(h(t_1)).
\]
Parameter gradients accumulate through
\[
\frac{dL}{d\theta}
=-\int_{t_1}^{t_0}
a(t)^\top \frac{\partial f_\theta}{\partial \theta}(h(t),t)\,dt,
\]
with signs depending on the chosen integration direction. The important idea is that gradients are themselves governed by a differential equation tied to the Jacobian of the dynamics.

\paragraph{Worked mechanism: ResNet as Euler discretization.}
Start from the ODE
\[
\dot h=f_\theta(h,t).
\]
Forward Euler with step \(\Delta t\) gives
\[
h_{k+1}=h_k+\Delta t\,f_\theta(h_k,t_k).
\]
If \(\Delta t=1\), this has the same algebraic form as a residual block
\[
h_{k+1}=h_k+F_{\theta_k}(h_k).
\]
The continuous-depth view asks what happens when many small residual blocks approximate a smooth flow. It also warns that changing depth, step size, or solver changes the effective computation.

\paragraph{Examples and connections.}
A simple continuous-depth model is
\[
\dot h=Wh,\qquad h(0)=x.
\]
Then
\[
h(1)=e^W x.
\]
This is a learned linear flow. If \(W\) has eigenvalues with large positive real parts, the hidden state can explode; if all real parts are negative, the flow contracts.

In computational neuroscience, latent neural dynamics models often use
\[
\dot z=f_\theta(z,u(t))
\]
to describe low-dimensional neural trajectories driven by inputs. Such models can interpolate between observations, but their biological interpretation depends on measurement assumptions and model validation.

In AI and scientific ML, continuous-depth models include neural ODEs, continuous normalizing flows, physics-informed neural networks, and neural operators. PDE-informed models may learn fields \(u(x,t)\) that satisfy approximate differential constraints, while neural ODEs learn finite-dimensional flows.

\paragraph{Common misconceptions and failure modes.}
A neural ODE is not automatically more accurate than a discrete network. Solver tolerances, stiffness, and model misspecification matter. The adjoint method is not magic memory-free differentiation without tradeoffs; numerical mismatch between forward and backward solves can affect gradients. Continuous depth also imposes structural constraints: an ODE flow with a unique solution cannot cross trajectories in state space, which may or may not suit the modeling task.

\paragraph{Exercises.}
\begin{enumerate}[label=\textbf{26.4.\arabic*.}, leftmargin=*]
\item \textbf{Mechanical.} Write the Euler update for \(\dot h=f_\theta(h,t)\) with step size \(\Delta t\).
\item \textbf{Proof or derivation.} Starting from \(h_{k+1}=h_k+\Delta t f_\theta(h_k,t_k)\), explain how a residual block corresponds to Euler with \(\Delta t=1\) and \(F_{\theta_k}=f_\theta(\cdot,t_k)\).
\item \textbf{Numerical/computational.} For \(\dot h=-2h\), \(h(0)=1\), compute Euler approximations to \(h(1)\) using \(\Delta t=0.5\) and \(\Delta t=0.25\). Compare qualitatively to \(e^{-2}\).
\item \textbf{Interpretation.} Why can solver tolerance affect both predictions and gradients in a neural ODE?
\item \textbf{Misconception/debugging.} A student says, "Continuous-depth models remove discretization error because they are continuous." Explain why any computer implementation still uses a numerical solver.
\end{enumerate}

\paragraph{Closing bridge.}
PDEs extend dynamics from trajectories to fields. They make boundary conditions, spatial operators, conservation, and discretization unavoidable. The next chapter returns to time evolution but adds randomness. In stochastic differential equations, symbols such as \(dW_t\) are not ordinary differentials; they represent random increments with scaling rules different from the deterministic \(dt\) used in ODEs and PDEs.

\clearpage
